\maketitle
\begin{abstract}
	Whether larval zebrafish can associate events separated by a stimulus-free interval remains unclear. Trace conditioning, in which a temporal gap separates the conditioned stimulus (CS) and unconditioned stimulus (US), provides a framework for studying how nervous systems link temporally discontinuous events. We developed a head-fixed classical conditioning assay in larval zebrafish that pairs a visual CS with an aversive optovin-mediated US. Larvae acquired both delay conditioning and trace conditioning across a 3-s stimulus-free interval. Learning was expressed as a reduction in tail-movement vigor that emerged during training, was absent in unpaired controls, became restricted to periods preceding the expected US, and extinguished during testing. Using comparable inter-trial intervals, conditioning was not detected at the population level with a 10-s trace interval. These findings establish predictive motor suppression as a behavioral readout of associative learning and provide a foundation for future cellular-resolution studies of neural mechanisms supporting associations across temporal gaps.
	\textbf{Keywords:} head-restrained behavior, larval zebrafish, delay conditioning, trace conditioning, optovin
\end{abstract}

%todo Excluded in this first minimal version: Learner classification; Temporal organization of learning; Temporal estimation / interval timing; 3-hour retention; MK-801; ca8 ablation; Imaging